\documentclass[10pt,a4paper]{amsart}
\usepackage{amsmath,amssymb,amsthm,graphicx}
% Retained preamble for the Stacks Project chapter snapshot in this directory.
% The chapter's own preamble.tex is a copy of the shared Stacks Project preamble; the block
% below is sliced MECHANICALLY, line for line, from it. Nothing was added, removed or
% redefined; the wrapper class, the xr-hyper/hyperref wiring and the externaldocument list
% are omitted because the assembler supplies the document class.
%
% lines taken from preamble.tex: 3, 5, 9, 11, 12, 156, 157, 158, ... (51 lines total)
%
% package coverage restored from the chapter's own preamble (2026-10-13): xy, xy-2cell

\usepackage[all]{xy}\usepackage{verbatim}
\xyoption{2cell}
\UseAllTwocells

\newenvironment{reference}{\comment}{\endcomment}
\newenvironment{slogan}{\comment}{\endcomment}
\newenvironment{history}{\comment}{\endcomment}
\theoremstyle{plain}
\newtheorem{theorem}[subsection]{Theorem}
\newtheorem{proposition}[subsection]{Proposition}
\newtheorem{lemma}[subsection]{Lemma}
\theoremstyle{definition}
\newtheorem{definition}[subsection]{Definition}
\newtheorem{example}[subsection]{Example}
\newtheorem{exercise}[subsection]{Exercise}
\newtheorem{situation}[subsection]{Situation}
\theoremstyle{remark}
\newtheorem{remark}[subsection]{Remark}
\newtheorem{remarks}[subsection]{Remarks}
\numberwithin{equation}{subsection}
\def\lim{\mathop{\mathrm{lim}}\nolimits}
\def\colim{\mathop{\mathrm{colim}}\nolimits}
\def\Spec{\mathop{\mathrm{Spec}}}
\def\Hom{\mathop{\mathrm{Hom}}\nolimits}
\def\Ext{\mathop{\mathrm{Ext}}\nolimits}
\def\SheafHom{\mathop{\mathcal{H}\!\mathit{om}}\nolimits}
\def\SheafExt{\mathop{\mathcal{E}\!\mathit{xt}}\nolimits}
\def\Sch{\mathit{Sch}}
\def\Mor{\mathop{\mathrm{Mor}}\nolimits}
\def\Ob{\mathop{\mathrm{Ob}}\nolimits}
\def\Sh{\mathop{\mathit{Sh}}\nolimits}
\def\NL{\mathop{N\!L}\nolimits}
\def\CH{\mathop{\mathrm{CH}}\nolimits}
\def\proetale{{pro\text{-}\acute{e}tale}}
\def\etale{{\acute{e}tale}}
\def\QCoh{\mathit{QCoh}}
\def\Ker{\mathop{\mathrm{Ker}}}
\def\Im{\mathop{\mathrm{Im}}}
\def\Coker{\mathop{\mathrm{Coker}}}
\def\Coim{\mathop{\mathrm{Coim}}}
\DeclareMathSymbol{\boxtimes}{\mathbin}{AMSa}{"02}
\def\QCohstack{\mathcal{QC}\!\mathit{oh}}
\def\Cohstack{\mathcal{C}\!\mathit{oh}}
\def\Spacesstack{\mathcal{S}\!\mathit{paces}}
\def\Quotfunctor{\mathrm{Quot}}
\def\Hilbfunctor{\mathrm{Hilb}}
\def\Curvesstack{\mathcal{C}\!\mathit{urves}}
\def\Polarizedstack{\mathcal{P}\!\mathit{olarized}}
\def\Complexesstack{\mathcal{C}\!\mathit{omplexes}}
\def\Pic{\mathop{\mathrm{Pic}}\nolimits}
\def\Picardstack{\mathcal{P}\!\mathit{ic}}
\def\Picardfunctor{\mathrm{Pic}}
\def\Deformationcategory{\mathcal{D}\!\mathit{ef}}
\usepackage{cleveref}
\crefname{originallemma}{original lemma}{original lemmas}
\Crefname{originallemma}{Original lemma}{Original lemmas}
\makeatletter
\newlabel{artin-lemma-ses-inf-and-T}{{0DNN}{source}{}{sourcepointer-1}{}}
\newlabel{artin-lemma-ses-inf-and-T@cref}{{[originallemma][0][]0DNN}{[1][0][]source}{}{}{}}
\makeatother
\crefname{originallemma}{original lemma}{original lemmas}
\Crefname{originallemma}{Original lemma}{Original lemmas}
\makeatletter
\newlabel{criteria-lemma-representable-by-spaces-formally-smooth}{{06D3}{source}{}{sourcepointer-2}{}}
\newlabel{criteria-lemma-representable-by-spaces-formally-smooth@cref}{{[originallemma][0][]06D3}{[1][0][]source}{}{}{}}
\makeatother
\crefname{originallemma}{original lemma}{original lemmas}
\Crefname{originallemma}{Original lemma}{Original lemmas}
\makeatletter
\newlabel{lemma-sheaf-of-infinitesimal-lifts}{{lemma-sheaf-of-infinitesimal-lifts}{source}{}{sourcepointer-3}{}}
\newlabel{lemma-sheaf-of-infinitesimal-lifts@cref}{{[originallemma][0][]lemma-sheaf-of-infinitesimal-lifts}{[1][0][]source}{}{}{}}
\makeatother
\crefname{originallemma}{original lemma}{original lemmas}
\Crefname{originallemma}{Original lemma}{Original lemmas}
\makeatletter
\newlabel{more-algebra-lemma-pseudo-coherence-and-base-change-ext}{{087R}{source}{}{sourcepointer-4}{}}
\newlabel{more-algebra-lemma-pseudo-coherence-and-base-change-ext@cref}{{[originallemma][0][]087R}{[1][0][]source}{}{}{}}
\makeatother
\crefname{originallemma}{original lemma}{original lemmas}
\Crefname{originallemma}{Original lemma}{Original lemmas}
\makeatletter
\newlabel{spaces-more-morphisms-lemma-finite-presentation-differentials}{{05ZF}{source}{}{sourcepointer-5}{}}
\newlabel{spaces-more-morphisms-lemma-finite-presentation-differentials@cref}{{[originallemma][0][]05ZF}{[1][0][]source}{}{}{}}
\makeatother
\crefname{originallemma}{original lemma}{original lemmas}
\Crefname{originallemma}{Original lemma}{Original lemmas}
\makeatletter
\newlabel{spaces-more-morphisms-lemma-smooth-formally-smooth}{{04AM}{source}{}{sourcepointer-6}{}}
\newlabel{spaces-more-morphisms-lemma-smooth-formally-smooth@cref}{{[originallemma][0][]04AM}{[1][0][]source}{}{}{}}
\makeatother
\crefname{originallemma}{original lemma}{original lemmas}
\Crefname{originallemma}{Original lemma}{Original lemmas}
\makeatletter
\newlabel{spaces-more-morphisms-section-action-by-derivations}{{0617}{source}{}{sourcepointer-7}{}}
\newlabel{spaces-more-morphisms-section-action-by-derivations@cref}{{[originallemma][0][]0617}{[1][0][]source}{}{}{}}
\makeatother
\crefname{originallemma}{original lemma}{original lemmas}
\Crefname{originallemma}{Original lemma}{Original lemmas}
\makeatletter
\newlabel{stacks-morphisms-lemma-finite-presentation-permanence}{{06Q6}{source}{}{sourcepointer-8}{}}
\newlabel{stacks-morphisms-lemma-finite-presentation-permanence@cref}{{[originallemma][0][]06Q6}{[1][0][]source}{}{}{}}
\makeatother
\crefname{originallemma}{original lemma}{original lemmas}
\Crefname{originallemma}{Original lemma}{Original lemmas}
\makeatletter
\newlabel{stacks-sheaves-lemma-surjective-flat-locally-finite-presentation}{{06XH}{source}{}{sourcepointer-9}{}}
\newlabel{stacks-sheaves-lemma-surjective-flat-locally-finite-presentation@cref}{{[originallemma][0][]06XH}{[1][0][]source}{}{}{}}
\makeatother
\title{Tangent spaces base-change along flat ring maps}
\author{The Stacks Project Authors}
\date{stacks-project:stacks-more-morphisms.tex:local-snapshot-20261012; GNU FDL 1.2 or later}
\begin{document}
\maketitle
\noindent\textit{Editorial note.} Statement span is lines 1698-1705 of sources/stacks-stacks-more-morphisms/pinned-original-source.tex; proof span is lines 1707-1793, the complete author proof (85 body lines).\par\medskip
\noindent\textit{Editorial note.} Navigation stubs: 9 cited-but-not-included labels are supplied as non-load-bearing pointers with their pinned tags where the snapshot tag index has them.\par\medskip
\noindent\textit{Editorial note.} extra\_packages is empty on purpose: the retained preamble already loads xy with options and a second declaration clashes in the assembler wrapper.\par\medskip
\noindent\textit{Editorial note.} Text is retained verbatim, including the author's own typos and placeholders.\par\medskip
\noindent\textit{Editorial note.} Queue index 9 of queue\_stacks-stacks-more-morphisms.json (pb=3601, stmt 1698-1705, proof 1707-1793).\par\medskip
\noindent\textit{Editorial note.} stacks-stacks-more-morphisms wave smm01 (2026-10-14): fresh adjudication accept per lane-s3/smm\_judgments.md; pool per Chief 10:5x open-pool ruling (net 59 with-proof rows, bloodline lane-s3 by anchor audit, v2 true-key namespace stacks:stacks-more-morphisms.tex, row lemma-deform-property-fp-over-ft avoided on cross-chapter key collision); 60 with-proof rows judged 45 accept / 14 reject / 1 avoid; patch-sd-envs contract spans scanned.\par\medskip
\noindent\textit{Allowed standard prerequisites (editorial).}\begin{itemize}
\item Explicit prerequisite boundary. The cited inputs are prerequisites; the argument in the retained span is not replaced by any cited result.
\end{itemize}\medskip
\ifdefined\hypertarget\hypertarget{sourcepointer-1}{}\fi
\noindent\textit{Source reference.} Original lemma 0DNN: sources/stacks-stacks-more-morphisms/pinned-original-source.tex, source lines 1793--1793. Cited by the retained statement or proof but not part of the collected spans; supplied only as a navigation pointer. The argument does not depend on its content.\par\medskip
\ifdefined\hypertarget\hypertarget{sourcepointer-2}{}\fi
\noindent\textit{Source reference.} Original lemma 06D3: sources/stacks-stacks-more-morphisms/pinned-original-source.tex, source lines 1793--1793. Cited by the retained statement or proof but not part of the collected spans; supplied only as a navigation pointer. The argument does not depend on its content.\par\medskip
\ifdefined\hypertarget\hypertarget{sourcepointer-3}{}\fi
\noindent\textit{Source reference.} Original lemma lemma-sheaf-of-infinitesimal-lifts: sources/stacks-stacks-more-morphisms/pinned-original-source.tex, source lines 1607--1618. Cited by the retained statement or proof but not part of the collected spans; supplied only as a navigation pointer. The argument does not depend on its content.\par\medskip
\ifdefined\hypertarget\hypertarget{sourcepointer-4}{}\fi
\noindent\textit{Source reference.} Original lemma 087R: sources/stacks-stacks-more-morphisms/pinned-original-source.tex, source lines 1793--1793. Cited by the retained statement or proof but not part of the collected spans; supplied only as a navigation pointer. The argument does not depend on its content.\par\medskip
\ifdefined\hypertarget\hypertarget{sourcepointer-5}{}\fi
\noindent\textit{Source reference.} Original lemma 05ZF: sources/stacks-stacks-more-morphisms/pinned-original-source.tex, source lines 1793--1793. Cited by the retained statement or proof but not part of the collected spans; supplied only as a navigation pointer. The argument does not depend on its content.\par\medskip
\ifdefined\hypertarget\hypertarget{sourcepointer-6}{}\fi
\noindent\textit{Source reference.} Original lemma 04AM: sources/stacks-stacks-more-morphisms/pinned-original-source.tex, source lines 1793--1793. Cited by the retained statement or proof but not part of the collected spans; supplied only as a navigation pointer. The argument does not depend on its content.\par\medskip
\ifdefined\hypertarget\hypertarget{sourcepointer-7}{}\fi
\noindent\textit{Source reference.} Original lemma 0617: sources/stacks-stacks-more-morphisms/pinned-original-source.tex, source lines 1793--1793. Cited by the retained statement or proof but not part of the collected spans; supplied only as a navigation pointer. The argument does not depend on its content.\par\medskip
\ifdefined\hypertarget\hypertarget{sourcepointer-8}{}\fi
\noindent\textit{Source reference.} Original lemma 06Q6: sources/stacks-stacks-more-morphisms/pinned-original-source.tex, source lines 1793--1793. Cited by the retained statement or proof but not part of the collected spans; supplied only as a navigation pointer. The argument does not depend on its content.\par\medskip
\ifdefined\hypertarget\hypertarget{sourcepointer-9}{}\fi
\noindent\textit{Source reference.} Original lemma 06XH: sources/stacks-stacks-more-morphisms/pinned-original-source.tex, source lines 1793--1793. Cited by the retained statement or proof but not part of the collected spans; supplied only as a navigation pointer. The argument does not depend on its content.\par\medskip
\begin{lemma}
\label{lemma-T-quasi-coherent}
Let $\mathcal{X}$ be an algebraic stack over a scheme $S$ whose
structure morphism $\mathcal{X} \to S$ is locally of finite presentation.
Let $A \to B$ be a flat $S$-algebra homomorphism.
Let $x$ be an object of $\mathcal{X}$ over $A$.
Then $T_x(M) \otimes_A B = T_y(M \otimes_A B)$.
\end{lemma}

\begin{proof}
Choose a scheme $U$ and a surjective smooth morphism $U \to \mathcal{X}$.
We first reduce the lemma to the case where $x$ lifts to $U$.
Recall that $T_x(M)$ is the set of isomorphism classes of lifts of $x$
to $A[M]$. Therefore
Lemma \ref{lemma-sheaf-of-infinitesimal-lifts}\footnote{This lemma applies:
$\Delta : \mathcal{X} \to \mathcal{X} \times_S \mathcal{X}$
is locally of finite presentation by
Morphisms of Stacks, Lemma
\ref{stacks-morphisms-lemma-finite-presentation-permanence}
and the assumption that $\mathcal{X} \to S$ is locally of finite presentation.
Therefore $\mathcal{I}_\mathcal{X} \to \mathcal{X}$ is locally
of finite presentation as a base change of $\Delta$.}
says that the rule
$$
A_1 \mapsto T_{x|_{A_1}}(M \otimes_A A_1)
$$
is a sheaf on the small \'etale site of $\Spec(A)$; the tensor product
is needed to make $A[M] \to A_1[M \otimes_A A_1]$ a flat ring map.
We may choose a faithfully flat \'etale ring map $A \to A_1$
such that $x|_{A_1}$ lifts to a morphism $u_1 : \Spec(A_1) \to U$, see
for example Sheaves on Stacks, Lemma
\ref{stacks-sheaves-lemma-surjective-flat-locally-finite-presentation}.
Write $A_2 = A_1 \otimes_A A_1$ and set $B_1 = B \otimes_A A_1$
and $B_2 = B \otimes_A A_2$. Consider the diagram
$$
\xymatrix{
0 \ar[r] &
T_y(M \otimes_A B) \ar[r] &
T_{y|_{B_1}}(M \otimes_A B_1) \ar[r] &
T_{y|_{B_2}}(M \otimes_A B_2) \\
0 \ar[r] &
T_x(M) \ar[r] \ar[u] &
T_{x|_{A_1}}(M \otimes_A A_1) \ar[r] \ar[u] &
T_{x|_{A_2}}(M \otimes_A A_2) \ar[u]
}
$$
The rows are exact by the sheaf condition. We have
$M \otimes_A B_i = (M \otimes_A A_i) \otimes_{A_i} B_i$.
Thus if we prove the result for the middle and right vertical arrow, then
the result follows. This reduces us to the case discussed in
the next paragraph.

\medskip\noindent
Assume that $x$ is the image of a morphism $u : \Spec(A) \to U$.
Observe that $T_u(M) \to T_x(M)$ is surjective since
$U \to \mathcal{X}$ is smooth and representable by
algebraic spaces, see Criteria for Representability, Lemma
\ref{criteria-lemma-representable-by-spaces-formally-smooth}
(see discussion preceding it for explanation) and
More on Morphisms of Spaces, Lemma
\ref{spaces-more-morphisms-lemma-smooth-formally-smooth}.
Set $R = U \times_\mathcal{X} U$. Recall that we obtain
a groupoid $(U, R, s, t, c, e, i)$ in algebraic spaces with
$\mathcal{X} = [U/R]$. By
Artin's Axioms, Lemma \ref{artin-lemma-ses-inf-and-T}
we have an exact sequence
$$
T_{e \circ u}(M) \to T_u(M) \oplus T_u(M) \to T_x(M) \to 0
$$
where the zero on the right was shown above. A similar sequence
holds for the base change to $B$. Thus the result we want follows if
we can prove the result of the lemma for
$T_u(M)$ and $T_{e \circ u}(M)$. This reduces us to the case discussed
in the next paragraph.

\medskip\noindent
Assume that $\mathcal{X} = X$ is an algebraic space locally of finite
presentation over $S$. Then we have
$$
T_x(M) = \Hom_A(x^*\Omega_{X/S}, M)
$$
by the discussion in More on Morphisms of Spaces, Section
\ref{spaces-more-morphisms-section-action-by-derivations}.
By the same token
$$
T_y(M \otimes_A B) = \Hom_B(y^*\Omega_{X/S}, M \otimes_A B)
$$
Since $X \to S$ is locally of finite presentation, we see that
$\Omega_{X/S}$ is locally of finite presentation, see
More on Morphisms of Spaces, Lemma
\ref{spaces-more-morphisms-lemma-finite-presentation-differentials}.
Hence $x^*\Omega_{X/S}$ is a finitely presented $A$-module.
Clearly, we have $y^*\Omega_{X/S} = x^*\Omega_{X/S} \otimes_A B$.
we conclude by More on Algebra, Lemma
\ref{more-algebra-lemma-pseudo-coherence-and-base-change-ext}.
\end{proof}

\end{document}
